We resort to resource augmentation analysis, introduced by of
Kalyanasundaram and Pruhs~\cite{KalyanasundaramP95}, to overcome the
above lower bounds. In this analysis the online algorithm is given a
faster machine than the optimal offline algorithm. For $s \ge 1$, an
algorithm $A$ is said to be $s$-speed $r$-competitive if $A$, when
given $s$-speed machine(s), achieves a competitive ratio of $r$
compared to the optimum  schedule given 1-speed machine(s).
We prove the following.
\begin{itemize}
\item For unicast setting, for any $\eps \in (0,1]$, there are
  $(1+\eps)$-speed $O(1/\eps)$-competitive algorithms in both single
  and multiple machine cases. Moreover, the algorithm for the multiple
  machine case immediately dispatches an arriving request to a
  machine, and is non-migratory. An algorithm is non-migratory
  if it processes each request on a single machine to which it is
  first assigned.

\item For broadcast setting, for any $\eps \in (0,1]$, there is a
  $(1+\eps)$-speed $O(1/\eps^2)$-competitive algorithm for uniform
  pages.  For non-uniform sized pages, for any $\eps \in (0,1]$, there
  is a $(1+\eps)$-speed $O(1/\eps^3)$-competitive algorithm.
\end{itemize}
